\section{Discusi\'on}
\label{sec:discusion}

\subsection{Interpretaci\'on de los resultados}

El mapeo realizado revela que \speckit{} no es, ni pretende ser, una herramienta de ingenier\'ia de requisitos conforme a IREB. Su cobertura de las actividades fundamentales es desigual: fuerte en documentaci\'on estructurada (5/10), d\'ebil en validaci\'on (2/10) y ausente en elicitaci\'on (1/10) y gesti\'on de requisitos (2/10). Esta distribuci\'on no es una deficiencia, sino una consecuencia de su prop\'osito: \speckit{} est\'a dise\~nado para la generaci\'on de c\'odigo, no para la gesti\'on del ciclo de vida de los requisitos.

Sin embargo, el an\'alisis revela que la brecha no es insalvable. La arquitectura de \speckit{} ---basada en plantillas, comandos extensibles, presets y \emph{workflows}--- ofrece puntos de entrada claros para incorporar pr\'acticas IREB. Un ``preset IREB'' podr\'ia modificar las plantillas de especificaci\'on para incluir atributos como tipo, fuente, \emph{rationale} y prioridad, mientras que una extensi\'on podr\'ia a\~nadir verificaci\'on de consistencia con criterios ISO 29148.

\subsection{Implicaciones para la pr\'actica}

Los resultados sugieren que \speckit{} puede utilizarse dentro de un flujo IREB siempre que la elicitaci\'on y la formalizaci\'on de requisitos se realicen previamente de forma manual o con otras herramientas. El pipeline SDD funciona como un ``generador de \'ultima milla'' que transforma requisitos ya formalizados en implementaciones, pero no debe esperarse que realice las tareas de descubrimiento y an\'alisis que corresponden al ingeniero de requisitos.

Esta separaci\'on de responsabilidades es coherente con la visi\'on de IREB: la herramienta de generaci\'on de c\'odigo no reemplaza al ingeniero de requisitos, sino que automatiza la parte mec\'anica de la transformaci\'on especificaci\'on--implementaci\'on.

\subsection{Limitaciones del estudio}

El presente trabajo tiene varias limitaciones que deben considerarse al interpretar sus resultados:

\begin{itemize}
    \item \textbf{Corpus limitado}: 52 requisitos de 6 proyectos, aunque representativos de distintos tipos de sistema, no constituyen una muestra estad\'isticamente significativa.
    \item \textbf{Fuentes de elicitaci\'on acotadas}: los \emph{changelogs} son la fuente principal. No se incluyeron sistem\'aticamente \emph{issues}, \emph{diffs} completos ni revisiones de c\'odigo como fuentes complementarias.
    \item \textbf{Ejecuci\'on de SpecKit pendiente}: los resultados del mapeo din\'amico (c\'omo responde \speckit{} a los requisitos) est\'an pendientes de la ejecuci\'on del pipeline.
    \item \textbf{Sin validaci\'on externa}: la r\'ubrica SRCI es aplicada por un \'unico evaluador. No hay inter- evaluador que permita calibrar la consistencia de las clasificaciones.
    \item \textbf{Alcance exploratorio}: el trabajo se enmarca como estudio de caso exploratorio, no como experimento controlado. Los patrones identificados son cualitativos y no generalizables estad\'isticamente.
\end{itemize}

\subsection{Trabajo futuro}

Las siguientes l\'ineas de trabajo podr\'ian extender los resultados obtenidos:

\begin{enumerate}
    \item \textbf{Preset IREB para SpecKit}: desarrollar y evaluar un \emph{preset} que modifique las plantillas de \speckit{} para incluir atributos IREB (tipo, fuente, \emph{rationale}, prioridad) y \emph{checklists} basados en ISO 29148.
    \item \textbf{Extensi\'on de elicitaci\'on}: implementar un comando \texttt{/speckit.elicit} que asista en la identificaci\'on de \emph{stakeholders}, fuentes y contexto del sistema.
    \item \textbf{Evaluaci\'on con m\'ultiples evaluadores}: replicar el an\'alisis SRCI con varios evaluadores para calibrar la consistencia de la r\'ubrica.
    \item \textbf{Ampliaci\'on del corpus}: incorporar proyectos de otros dominios (sistemas cr\'iticos, IoT, tiempo real) para evaluar la generalidad del mapeo.
\end{enumerate}
